\documentclass[12pt,a4paper]{article}

\usepackage[utf8]{inputenc}
\usepackage[T5]{fontenc}
\usepackage[vietnamese]{babel}
\usepackage{geometry}
\usepackage{graphicx}
\usepackage{float}
\usepackage{array}
\usepackage{longtable}
\usepackage{hyperref}
\usepackage{verbatim}

\geometry{
    left=3cm,
    right=2cm,
    top=2.5cm,
    bottom=2.5cm
}

\renewcommand{\arraystretch}{1.3}
\setlength{\parindent}{1cm}
\setlength{\parskip}{6pt}

\newcommand{\safeimage}[3]{
    \begin{figure}[H]
        \centering
        \IfFileExists{#1}{
            \includegraphics[width=#2\textwidth]{#1}
        }{
            \fbox{
                \begin{minipage}[c][5cm][c]{#2\textwidth}
                    \centering
                    Chưa có ảnh: \texttt{#1}
                \end{minipage}
            }
        }
        \caption{#3}
    \end{figure}
}

\begin{document}

% ================= TRANG BÌA =================
\begin{titlepage}
    \centering

    {\Large TRƯỜNG ĐẠI HỌC ...}\\
    {\Large KHOA ...}\\[2cm]

    {\LARGE \textbf{BÁO CÁO BÀI TẬP LỚN}}\\[0.5cm]
    {\Large \textbf{MÔN: LẬP TRÌNH HƯỚNG ĐỐI TƯỢNG}}\\[1.5cm]

    {\Huge \textbf{WILD-LIFE ECO SIMULATION}}\\[0.5cm]
    {\Large \textbf{Hệ thống mô phỏng hệ sinh thái hoang dã}}\\[2cm]

    \begin{flushleft}
        \textbf{Giảng viên hướng dẫn:} ...\\
        \textbf{Nhóm thực hiện:} ...\\
        \textbf{Lớp:} ...\\
    \end{flushleft}

    \vspace{1cm}

    \begin{tabular}{|c|l|c|}
        \hline
        \textbf{STT} & \textbf{Họ và tên} & \textbf{MSSV} \\
        \hline
        1 & TungTran / TungSoYeu & ... \\
        \hline
        2 & davidtaiduc & ... \\
        \hline
        3 & Thành viên khác nếu có & ... \\
        \hline
    \end{tabular}

    \vfill
    {\large TP. Hồ Chí Minh, tháng 06 năm 2026}
\end{titlepage}

\tableofcontents
\newpage

% ================= NỘI DUNG =================

\section{Giới thiệu đề tài}

\subsection{Bối cảnh}

Đề tài \textbf{Wild-Life Eco Simulation} là một ứng dụng JavaFX mô phỏng hệ sinh thái hoang dã theo hướng lập trình hướng đối tượng. Chương trình xây dựng một bản đồ hai chiều gồm nhiều loại địa hình như đồng cỏ, rừng rậm, hồ nước, bùn, vách đá và bụi rậm. Trên bản đồ có nhiều nhóm thực thể khác nhau, bao gồm thực vật, động vật ăn cỏ, động vật săn mồi, sinh vật sống dưới nước và thợ săn.

Trong mô phỏng, mỗi thực thể không chỉ được vẽ lên màn hình mà còn có trạng thái riêng. Ví dụ, động vật có máu, độ đói, độ khát, tốc độ, tầm nhìn, con mồi, kẻ thù tự nhiên và chiến lược sinh tồn. Thực vật có khả năng phát triển, bị ăn và lan rộng. Các mùa trong năm cũng ảnh hưởng đến tốc độ phát triển của thực vật và khả năng sinh sản của động vật.

\subsection{Mục tiêu}

Mục tiêu chính của đề tài là áp dụng các kỹ thuật lập trình hướng đối tượng để xây dựng một hệ thống mô phỏng có cấu trúc rõ ràng, dễ mở rộng và dễ kiểm thử. Các mục tiêu cụ thể gồm:

\begin{itemize}
    \item Xây dựng hệ thống lớp thể hiện được quan hệ kế thừa giữa thực thể, động vật và thực vật.
    \item Tách biệt logic mô phỏng, logic hiển thị và logic hành vi.
    \item Áp dụng các kỹ thuật OOP như đóng gói, kế thừa, đa hình, trừu tượng hóa, interface và design pattern.
    \item Cung cấp giao diện trực quan để người dùng quan sát, điều khiển và thêm thực thể vào môi trường mô phỏng.
    \item Có dữ liệu mô phỏng ban đầu, thống kê số lượng thực thể và một số test kiểm thử hành vi quan trọng.
\end{itemize}

\section{Phân công công việc các thành viên trong nhóm}

\begin{table}[H]
    \centering
    \caption{Phân công công việc và tỉ lệ đóng góp}
    \begin{tabular}{|c|p{4cm}|p{7cm}|c|}
        \hline
        \textbf{STT} & \textbf{Thành viên} & \textbf{Công việc phụ trách} & \textbf{Đóng góp} \\
        \hline
        1 & TungTran / TungSoYeu
        & Thiết kế cấu trúc project Maven, xây dựng engine mô phỏng, hệ thống model, logic mùa, sinh sản, di chuyển, va chạm và tài liệu hướng dẫn.
        & 55\% \\
        \hline
        2 & davidtaiduc
        & Hỗ trợ phát triển tính năng, chỉnh sửa và tích hợp một phần logic chương trình, hỗ trợ kiểm tra hoạt động sau khi ghép code.
        & 20\% \\
        \hline
        3 & Thành viên khác nếu có
        & Hỗ trợ kiểm thử, vẽ UML, chuẩn bị báo cáo, chụp ảnh demo chương trình và chuẩn bị nội dung thuyết trình.
        & 25\% \\
        \hline
    \end{tabular}
\end{table}

\subsection{Mô tả đóng góp cụ thể}

\begin{itemize}
    \item \textbf{Thiết kế kiến trúc:} nhóm chia chương trình thành các package \texttt{model}, \texttt{engine}, \texttt{strategy}, \texttt{view}, \texttt{sound} và \texttt{util}. Cách chia này giúp mỗi package đảm nhiệm một vai trò riêng.
    \item \textbf{Xây dựng mô hình dữ liệu:} nhóm cài đặt các lớp cha trừu tượng \texttt{Entity}, \texttt{Animal}, \texttt{Plant} và các lớp con cụ thể như \texttt{Rabbit}, \texttt{Deer}, \texttt{Wolf}, \texttt{Tiger}, \texttt{Grass}, \texttt{FruitTree}.
    \item \textbf{Xây dựng engine mô phỏng:} nhóm cài đặt \texttt{SimulationEngine}, \texttt{EntityManager} và \texttt{SeasonManager} để điều khiển vòng lặp mô phỏng, sinh sản, duy trì quần thể và dọn dẹp thực thể đã chết.
    \item \textbf{Xây dựng giao diện:} nhóm cài đặt \texttt{GameView}, \texttt{Renderer}, \texttt{BasicRenderer}, \texttt{SpriteRenderer}, \texttt{Camera} để hiển thị bản đồ, thực thể, trạng thái và bảng thống kê.
    \item \textbf{Kiểm thử:} nhóm viết các test JUnit để kiểm tra chiến lược săn mồi, đổi chiến lược khi đói, khả năng ẩn nấp trong bụi rậm và sát thương phục kích của hổ.
\end{itemize}

\section{Chi tiết nhật ký làm việc nhóm}

\begin{longtable}{|c|c|p{4cm}|p{6cm}|p{3cm}|}
    \hline
    \textbf{STT} & \textbf{Thời gian} & \textbf{Thành viên} & \textbf{Nội dung công việc} & \textbf{Kết quả} \\
    \hline
    1 & Giai đoạn 1 & Cả nhóm & Thảo luận đề tài, xác định phạm vi mô phỏng, các loại thực thể và yêu cầu cần thể hiện về OOP. & Thống nhất đề tài Wild-Life Eco Simulation. \\
    \hline
    2 & Giai đoạn 2 & TungTran / TungSoYeu & Khởi tạo project Java Maven, cấu hình JavaFX, JUnit, Maven Wrapper và module \texttt{com.ecosim}. & Project có thể build và chạy bằng Maven Wrapper. \\
    \hline
    3 & Giai đoạn 3 & Cả nhóm & Thiết kế package, xác định quan hệ giữa model, engine, strategy và view. & Có cấu trúc package rõ ràng, tách biệt trách nhiệm. \\
    \hline
    4 & Giai đoạn 4 & Cả nhóm & Cài đặt hệ thống thực thể, địa hình, mùa, bản đồ và dữ liệu khởi tạo ban đầu. & Hoàn thành phần model chính của hệ sinh thái. \\
    \hline
    5 & Giai đoạn 5 & Cả nhóm & Cài đặt vòng lặp mô phỏng, xử lý hành động, di chuyển, uống nước, ăn, tấn công và va chạm mềm. & Mô phỏng có thể chạy theo tick và thay đổi theo thời gian. \\
    \hline
    6 & Giai đoạn 6 & Cả nhóm & Cài đặt Strategy Pattern cho hành vi sinh tồn: thụ động, sợ hãi, săn mồi và liều lĩnh. & Động vật có hành vi khác nhau và có thể đổi chiến lược khi runtime. \\
    \hline
    7 & Giai đoạn 7 & Cả nhóm & Xây dựng giao diện JavaFX, canvas, camera, renderer, sprite, âm thanh và bảng thống kê. & Người dùng có thể quan sát và tương tác với mô phỏng. \\
    \hline
    8 & Giai đoạn 8 & Cả nhóm & Viết test, sửa lỗi, hoàn thiện README, vẽ UML, chuẩn bị báo cáo và ảnh demo. & Hoàn thiện phiên bản báo cáo và chương trình demo. \\
    \hline
\end{longtable}

\section{Thống kê dữ liệu thu thập được}

Đề tài này không thu thập dữ liệu khảo sát bên ngoài mà sử dụng dữ liệu cấu hình và dữ liệu mô phỏng được khởi tạo trong chương trình. Dữ liệu này giúp tạo môi trường ban đầu để kiểm tra sự cân bằng giữa thực vật, động vật ăn cỏ, động vật săn mồi và các đối tượng đặc biệt.

\begin{table}[H]
    \centering
    \caption{Thống kê dữ liệu mô phỏng}
    \begin{tabular}{|c|p{5cm}|c|p{5cm}|}
        \hline
        \textbf{STT} & \textbf{Loại dữ liệu} & \textbf{Số lượng} & \textbf{Nguồn dữ liệu} \\
        \hline
        1 & Kích thước bản đồ & 120 x 80 = 9600 tile & \texttt{Constants.MAP\_WIDTH}, \texttt{Constants.MAP\_HEIGHT} \\
        \hline
        2 & Loại địa hình & 6 & \texttt{TerrainType}: đồng cỏ, rừng, nước, bùn, đá, bụi rậm \\
        \hline
        3 & Mùa trong năm & 4 & \texttt{Season}: xuân, hạ, thu, đông \\
        \hline
        4 & Thực vật khởi tạo ban đầu & 150 & 120 cỏ, 30 cây ăn quả \\
        \hline
        5 & Động vật khởi tạo ban đầu & 78 & Thỏ, hươu, sói, hổ, voi, thợ săn, cá, vịt \\
        \hline
        6 & Tổng thực thể ban đầu & 228 & Tổng số thực vật và động vật khi bắt đầu mô phỏng \\
        \hline
        7 & Ảnh sprite PNG & 9 & Thư mục \texttt{src/main/resources/sprites} \\
        \hline
        8 & Tệp âm thanh WAV & 6 & Thư mục \texttt{src/main/resources/sounds} \\
        \hline
        9 & Test case JUnit & 5 & Các file test trong \texttt{src/test/java} \\
        \hline
    \end{tabular}
\end{table}

\begin{table}[H]
    \centering
    \caption{Số lượng thực thể khởi tạo ban đầu}
    \begin{tabular}{|c|l|c|p{6cm}|}
        \hline
        \textbf{STT} & \textbf{Loại thực thể} & \textbf{Số lượng} & \textbf{Vai trò trong mô phỏng} \\
        \hline
        1 & Cỏ & 120 & Nguồn thức ăn chính cho thỏ, hươu và voi. \\
        \hline
        2 & Cây ăn quả & 30 & Nguồn thức ăn trong vùng rừng, hỗ trợ động vật ăn thực vật. \\
        \hline
        3 & Thỏ & 30 & Động vật ăn cỏ, là con mồi của sói, hổ và thợ săn. \\
        \hline
        4 & Hươu & 15 & Động vật ăn thực vật, có khả năng chạy trốn khi gặp nguy hiểm. \\
        \hline
        5 & Sói & 5 & Động vật săn mồi, thường săn thỏ và hươu. \\
        \hline
        6 & Hổ & 2 & Động vật săn mồi mạnh hơn, có lợi thế khi ở địa hình rừng. \\
        \hline
        7 & Thợ săn & 1 & Thực thể đặc biệt, có tầm nhìn xa và săn nhiều loại con mồi. \\
        \hline
        8 & Voi & 2 & Động vật ưu tiên cao, có kích thước lớn và ít bị đe dọa. \\
        \hline
        9 & Cá & 15 & Sinh vật sống chủ yếu ở nước hoặc bùn. \\
        \hline
        10 & Vịt & 8 & Sinh vật di chuyển ở nước hoặc bùn, có hành vi thụ động. \\
        \hline
    \end{tabular}
\end{table}

\section{Biểu đồ UML}

\subsection{Biểu đồ phụ thuộc gói}

\safeimage{uml/package-diagram.png}{0.95}{Biểu đồ phụ thuộc gói của project}

Biểu đồ phụ thuộc gói mô tả cách các package trong project liên hệ với nhau. Package \texttt{com.ecosim} đóng vai trò khởi động ứng dụng thông qua lớp \texttt{App}. Từ đây, chương trình tạo \texttt{SimulationEngine} để xử lý logic mô phỏng và tạo \texttt{GameView} để hiển thị giao diện người dùng.

Package \texttt{engine} phụ thuộc vào \texttt{model} vì engine cần cập nhật trạng thái của bản đồ, mùa và các thực thể. Package này cũng phụ thuộc vào \texttt{strategy} vì mỗi động vật sẽ dùng một chiến lược sinh tồn để quyết định hành động tiếp theo. Package \texttt{view} phụ thuộc vào \texttt{engine} để điều khiển mô phỏng, phụ thuộc vào \texttt{model} để lấy dữ liệu cần vẽ, và phụ thuộc vào \texttt{sound} để phát âm thanh theo sự kiện. Package \texttt{util} được nhiều package khác sử dụng vì chứa hằng số cấu hình và lớp vector dùng chung.

Thiết kế theo package như trên giúp giảm sự phụ thuộc vòng và làm rõ trách nhiệm của từng phần. Khi cần thêm một loài mới, nhóm chủ yếu chỉnh trong \texttt{model} và có thể bổ sung hành vi trong \texttt{strategy}; khi cần đổi cách hiển thị, nhóm chỉ cần chỉnh trong \texttt{view} mà không ảnh hưởng lớn đến logic mô phỏng.

\subsection{Biểu đồ lớp}

\safeimage{uml/class-diagram.png}{0.95}{Biểu đồ lớp chính của project}

Biểu đồ lớp thể hiện cấu trúc kế thừa và quan hệ sử dụng giữa các lớp chính. Ở trung tâm của mô hình là lớp trừu tượng \texttt{Entity}, đại diện cho mọi đối tượng xuất hiện trong hệ sinh thái. Từ \texttt{Entity}, chương trình tách thành hai nhánh lớn là \texttt{Animal} và \texttt{Plant}. Nhánh \texttt{Animal} chứa các loài động vật như thỏ, hươu, sói, hổ, voi, thợ săn, cá và vịt. Nhánh \texttt{Plant} chứa cỏ và cây ăn quả.

Lớp \texttt{Animal} không tự quyết định toàn bộ hành vi mà ủy quyền cho interface \texttt{SurvivalStrategy}. Các lớp \texttt{PassiveStrategy}, \texttt{ScaredStrategy}, \texttt{HunterStrategy} và \texttt{AggressiveStrategy} là các cách cài đặt cụ thể của interface này. Nhờ vậy, hành vi của động vật có thể thay đổi linh hoạt mà không cần sửa trực tiếp code của từng loài.

Phần điều khiển mô phỏng nằm ở các lớp \texttt{SimulationEngine}, \texttt{EntityManager} và \texttt{SeasonManager}. \texttt{SimulationEngine} điều phối toàn bộ vòng lặp; \texttt{EntityManager} quản lý danh sách thực thể, xử lý sinh sản và va chạm; \texttt{SeasonManager} quản lý mùa hiện tại. Phần hiển thị được tách qua interface \texttt{Renderer}, với hai lớp cài đặt là \texttt{BasicRenderer} và \texttt{SpriteRenderer}. Cách tách này làm cho chương trình vừa có thể chạy logic mô phỏng, vừa có thể thay đổi giao diện hiển thị một cách độc lập.

\section{Giải thích thiết kế}

\subsection{Các gói}

\safeimage{uml/package-diagram.png}{0.95}{Sơ đồ các package trong project}

Sơ đồ package cho thấy project được tổ chức theo hướng tách trách nhiệm. Mỗi package đảm nhiệm một nhóm chức năng riêng, hạn chế việc đặt quá nhiều logic vào một lớp hoặc một thư mục duy nhất. Cách tổ chức này phù hợp với lập trình hướng đối tượng vì mỗi phần có vai trò rõ ràng và có thể được mở rộng độc lập.

\begin{table}[H]
    \centering
    \caption{Mô tả các package trong project}
    \begin{tabular}{|c|p{4cm}|p{8cm}|}
        \hline
        \textbf{STT} & \textbf{Package} & \textbf{Vai trò} \\
        \hline
        1 & \texttt{com.ecosim} & Chứa lớp khởi động ứng dụng \texttt{App}. Lớp này tạo engine, tạo giao diện, cấu hình cửa sổ và khởi chạy JavaFX. \\
        \hline
        2 & \texttt{com.ecosim.model} & Chứa các lớp mô hình chính như \texttt{Entity}, \texttt{Animal}, \texttt{Plant}, các loài động vật, thực vật, bản đồ, địa hình, mùa và hành động. \\
        \hline
        3 & \texttt{com.ecosim.engine} & Chứa logic điều khiển mô phỏng, bao gồm vòng lặp tick, quản lý thực thể, quản lý mùa, sinh sản, uống nước, duy trì quần thể và xử lý va chạm. \\
        \hline
        4 & \texttt{com.ecosim.strategy} & Chứa các chiến lược hành vi của động vật. Package này giúp tách thuật toán quyết định hành động khỏi lớp động vật. \\
        \hline
        5 & \texttt{com.ecosim.view} & Chứa giao diện JavaFX, canvas, camera, renderer, sprite và hiệu ứng. Đây là phần chịu trách nhiệm hiển thị và nhận tương tác người dùng. \\
        \hline
        6 & \texttt{com.ecosim.sound} & Chứa \texttt{SoundManager}, dùng để quản lý và phát âm thanh theo sự kiện trong mô phỏng. \\
        \hline
        7 & \texttt{com.ecosim.util} & Chứa các tiện ích dùng chung như \texttt{Constants} và \texttt{Vector2D}. Các lớp này được nhiều package khác sử dụng. \\
        \hline
    \end{tabular}
\end{table}

\subsection{Các lớp chính}

\safeimage{uml/class-diagram.png}{0.95}{Sơ đồ các lớp chính trong project}

Sơ đồ lớp được thiết kế theo hướng đặt các lớp trừu tượng ở phía trên và các lớp cụ thể ở phía dưới. Cấu trúc này giúp thể hiện rõ quan hệ kế thừa: \texttt{Entity} là lớp gốc, \texttt{Animal} và \texttt{Plant} là hai nhóm chính, còn các loài cụ thể là lớp con. Ngoài quan hệ kế thừa, sơ đồ còn thể hiện quan hệ sử dụng giữa engine, strategy, renderer và bản đồ.

\begin{longtable}{|c|p{4cm}|p{8cm}|}
    \hline
    \textbf{STT} & \textbf{Lớp} & \textbf{Mô tả chi tiết} \\
    \hline
    1 & \texttt{App} & Entry point của JavaFX. Lớp này tạo \texttt{SimulationEngine}, tạo \texttt{GameView}, cấu hình \texttt{Scene}, load CSS và hiển thị cửa sổ chương trình. \\
    \hline
    2 & \texttt{SimulationEngine} & Điều phối mỗi tick của mô phỏng. Lớp này cập nhật mùa, duy trì quần thể, cập nhật từng thực thể, gọi strategy để lấy hành động và xử lý các bước sau cùng như va chạm, uống nước, đổi strategy và cleanup. \\
    \hline
    3 & \texttt{EntityManager} & Quản lý danh sách thực thể. Lớp này chịu trách nhiệm spawn ban đầu, thêm thực thể mới, tìm entity gần trong tầm nhìn, sinh sản mùa xuân, duy trì số lượng tối thiểu và thống kê số lượng theo loại. \\
    \hline
    4 & \texttt{SeasonManager} & Quản lý mùa hiện tại, thời gian trôi qua trong mùa và tiến trình mùa. Lớp này giúp mô phỏng thay đổi theo xuân, hạ, thu, đông. \\
    \hline
    5 & \texttt{WorldMap} & Quản lý bản đồ dạng lưới tile. Lớp này lưu các ô địa hình, cung cấp hàm lấy địa hình tại tọa độ, đặt địa hình và tìm vị trí ngẫu nhiên theo loại terrain. \\
    \hline
    6 & \texttt{Entity} & Lớp cha trừu tượng cho mọi thực thể. Lớp này chứa id, tên, vị trí, trạng thái sống, độ ưu tiên và kích thước. \\
    \hline
    7 & \texttt{Animal} & Lớp cha trừu tượng cho động vật. Lớp này quản lý máu, độ đói, độ khát, tốc độ, hướng di chuyển, trạng thái, con mồi, kẻ thù và strategy hiện tại. \\
    \hline
    8 & \texttt{Plant} & Lớp cha trừu tượng cho thực vật. Lớp này quản lý quá trình tăng trưởng, khả năng bị ăn, lan rộng và tạo cây con. \\
    \hline
    9 & \texttt{Rabbit}, \texttt{Deer} & Nhóm động vật ăn thực vật. Chúng có kẻ thù tự nhiên và thường dùng \texttt{ScaredStrategy}; khi quá đói có thể chuyển sang \texttt{AggressiveStrategy}. \\
    \hline
    10 & \texttt{Wolf}, \texttt{Tiger}, \texttt{Hunter} & Nhóm săn mồi. Các lớp này dùng \texttt{HunterStrategy} để tìm con mồi trong tầm nhìn và di chuyển đến mục tiêu. \\
    \hline
    11 & \texttt{Elephant} & Động vật lớn, có độ ưu tiên cao khi xử lý va chạm. Voi chủ yếu ăn thực vật và gần như không có kẻ thù tự nhiên trong mô phỏng hiện tại. \\
    \hline
    12 & \texttt{Fish}, \texttt{Duck} & Nhóm sinh vật sống ở môi trường nước hoặc bùn, thể hiện việc mỗi loài có thể có quy tắc di chuyển theo địa hình khác nhau. \\
    \hline
    13 & \texttt{SurvivalStrategy} & Interface định nghĩa phương thức \texttt{decide()}, dùng để quyết định hành động kế tiếp của động vật dựa trên bản thân, các thực thể xung quanh và bản đồ. \\
    \hline
    14 & \texttt{Renderer} & Interface cho phần hiển thị. Hai lớp \texttt{BasicRenderer} và \texttt{SpriteRenderer} cài đặt interface này để hỗ trợ hai chế độ vẽ khác nhau. \\
    \hline
    15 & \texttt{GameView} & Lớp giao diện chính, chứa canvas, thanh công cụ, sidebar thống kê, menu thêm thực thể, xử lý chuột, zoom, pan và game loop. \\
    \hline
    16 & \texttt{AnimalReproductionFactory} & Lớp tiện ích tạo con non theo từng loài. Factory giúp gom logic sinh sản vào một nơi thay vì rải rác trong từng lớp động vật. \\
    \hline
\end{longtable}

\section{Các kỹ thuật lập trình hướng đối tượng đã áp dụng}

\begin{longtable}{|c|p{3.2cm}|p{5cm}|p{5cm}|}
    \hline
    \textbf{STT} & \textbf{Kỹ thuật} & \textbf{Áp dụng ở đâu} & \textbf{Lợi ích} \\
    \hline
    1 & Đóng gói & Các thuộc tính trong \texttt{Entity}, \texttt{Animal}, \texttt{WorldMap}, \texttt{EntityManager} được đặt \texttt{private} hoặc \texttt{protected}. Một số giá trị như hunger và thirst được cập nhật qua setter để giới hạn trong khoảng hợp lệ. & Bảo vệ dữ liệu nội bộ, tránh sửa trực tiếp gây lỗi và giúp kiểm soát trạng thái của đối tượng. \\
    \hline
    2 & Kế thừa & \texttt{Animal} và \texttt{Plant} kế thừa \texttt{Entity}; các loài cụ thể kế thừa \texttt{Animal} hoặc \texttt{Plant}. & Tái sử dụng code chung về vị trí, trạng thái sống, di chuyển, cập nhật và xử lý hành động. \\
    \hline
    3 & Đa hình & \texttt{SimulationEngine} xử lý danh sách \texttt{Entity}; \texttt{Renderer} có hai cài đặt là \texttt{BasicRenderer} và \texttt{SpriteRenderer}; \texttt{SurvivalStrategy} có nhiều lớp triển khai. & Cho phép chương trình xử lý nhiều đối tượng khác nhau thông qua cùng một kiểu cha hoặc interface. \\
    \hline
    4 & Trừu tượng hóa & \texttt{Entity}, \texttt{Animal}, \texttt{Plant} là abstract class; \texttt{Renderer} và \texttt{SurvivalStrategy} là interface. & Che giấu chi tiết cài đặt, tập trung vào hành vi chung cần có của từng nhóm đối tượng. \\
    \hline
    5 & Strategy Pattern & Các lớp \texttt{PassiveStrategy}, \texttt{ScaredStrategy}, \texttt{HunterStrategy}, \texttt{AggressiveStrategy} cùng cài đặt \texttt{SurvivalStrategy}. & Tách logic quyết định hành vi khỏi lớp động vật, giúp đổi chiến lược khi runtime và dễ thêm chiến lược mới. \\
    \hline
    6 & Factory & \texttt{AnimalReproductionFactory} tạo con non theo từng loài. & Gom logic sinh sản vào một lớp riêng, giúp code dễ bảo trì và dễ mở rộng. \\
    \hline
    7 & Tách trách nhiệm & \texttt{model} chứa dữ liệu và logic sinh học, \texttt{engine} điều khiển mô phỏng, \texttt{view} hiển thị giao diện, \texttt{strategy} quyết định hành vi. & Làm code rõ ràng, giảm phụ thuộc lẫn nhau và thuận tiện khi kiểm thử. \\
    \hline
    8 & Singleton đơn giản & \texttt{AssetManager.getInstance()} quản lý cache sprite. & Tránh tải lặp tài nguyên hình ảnh, giúp renderer sử dụng lại sprite hiệu quả hơn. \\
    \hline
\end{longtable}

\section{Công nghệ sử dụng và thuật toán}

\subsection{Công nghệ sử dụng}

\begin{itemize}
    \item \textbf{Ngôn ngữ lập trình:} Java 21.
    \item \textbf{Framework giao diện:} JavaFX 21, gồm JavaFX Controls, JavaFX FXML và JavaFX Media.
    \item \textbf{Công cụ build:} Maven, Maven Wrapper và JavaFX Maven Plugin.
    \item \textbf{Kiểm thử:} JUnit Jupiter 5.10.2 và Maven Surefire Plugin.
    \item \textbf{Tài nguyên giao diện:} CSS, ảnh sprite PNG và âm thanh WAV.
    \item \textbf{Công cụ báo cáo:} LaTeX và biểu đồ UML dạng ảnh đen trắng.
\end{itemize}

\subsection{Thuật toán và kỹ thuật xử lý nổi bật}

\begin{itemize}
    \item \textbf{Vòng lặp simulation theo tick:} \texttt{SimulationEngine.tick()} cập nhật toàn bộ mô phỏng theo \texttt{deltaTime}. Cách làm này giúp mô phỏng không phụ thuộc cứng vào số frame.
    \item \textbf{Tìm thực thể gần trong tầm nhìn:} \texttt{EntityManager.getNearby()} duyệt danh sách entity và lọc các đối tượng còn sống nằm trong bán kính quan sát của động vật.
    \item \textbf{Quyết định hành vi bằng strategy:} mỗi động vật đưa thông tin bản thân, danh sách đối tượng gần đó và bản đồ cho strategy. Strategy trả về một \texttt{Action} như đi lang thang, chạy trốn, ăn, uống nước hoặc tấn công.
    \item \textbf{Di chuyển có gia tốc và giảm tốc:} \texttt{Animal.doMoveTo()} tính tốc độ theo địa hình, giảm tốc khi gần mục tiêu và dùng vector để tạo chuyển động mượt hơn.
    \item \textbf{Xử lý va chạm mềm:} \texttt{EntityManager.resolveMovementPriority()} dùng độ ưu tiên để đẩy nhẹ thực thể ưu tiên thấp ra xa thực thể ưu tiên cao, tránh hiện tượng chồng lấn.
    \item \textbf{Sinh sản theo mùa:} khi bắt đầu mùa xuân, \texttt{EntityManager.processSpringReproduction()} xét điều kiện sinh sản, xác suất sinh sản và giới hạn quần thể.
    \item \textbf{Cache tài nguyên:} \texttt{AssetManager} lưu sprite đã tải để renderer không phải đọc lại ảnh nhiều lần.
\end{itemize}

\section{Hướng dẫn sử dụng chương trình}

\subsection{Cách chạy chương trình}

\begin{enumerate}
    \item Cài đặt Java 21 trở lên.
    \item Mở thư mục project trong VS Code, IntelliJ IDEA, Eclipse hoặc terminal.
    \item Chạy nhanh bằng lệnh:
\begin{verbatim}
.\run.bat
\end{verbatim}
    \item Hoặc chạy trực tiếp bằng Maven Wrapper:
\begin{verbatim}
.\mvnw.cmd clean javafx:run
\end{verbatim}
    \item Để build và chạy test:
\begin{verbatim}
.\mvnw.cmd package
.\mvnw.cmd test
\end{verbatim}
\end{enumerate}

\subsection{Các thao tác chính}

\begin{itemize}
    \item Nút \textbf{Bắt đầu/Tạm dừng}: chạy hoặc dừng mô phỏng.
    \item Mục \textbf{Tốc độ}: thay đổi tốc độ mô phỏng từ 0.5x đến 8x.
    \item Nút \textbf{Đồ họa}: đổi giữa chế độ Basic Shapes và Sprites.
    \item Mục \textbf{Xem vùng}: chuyển nhanh camera đến toàn bản đồ, đồng cỏ, rừng rậm hoặc hồ nước.
    \item Click trái với công cụ \textbf{Chọn}: chọn entity để xem thông tin chi tiết như máu, độ đói, độ khát, trạng thái và strategy.
    \item Click trái với công cụ \textbf{Gieo cỏ}, \textbf{Đặt đá}, \textbf{Đặt bụi rậm}: thay đổi môi trường mô phỏng.
    \item Click phải trên bản đồ: mở menu thêm thỏ, hươu, sói, hổ, voi, thợ săn, cá, vịt, cỏ hoặc cây ăn quả.
    \item Scroll chuột để zoom; kéo chuột giữa hoặc dùng công cụ \textbf{Chọn} để pan camera.
\end{itemize}

\subsection{Một số ảnh demo chương trình}

\safeimage{images/demo-main.png}{0.85}{Màn hình chính của chương trình}

Ảnh này nên thể hiện toàn bộ giao diện chính, bao gồm bản đồ mô phỏng, thanh điều khiển, sidebar thống kê và các thực thể đang di chuyển trên bản đồ.

\safeimage{images/demo-sprite-mode.png}{0.85}{Chế độ hiển thị sprite}

Ảnh này nên chụp lúc bật chế độ sprite để thể hiện sự khác biệt giữa hiển thị bằng hình cơ bản và hiển thị bằng ảnh nhân vật.

\safeimage{images/demo-selected-entity.png}{0.85}{Bảng thông tin khi chọn một entity}

Ảnh này nên chụp lúc chọn một động vật cụ thể. Sidebar cần hiển thị các thông tin như trạng thái, máu, độ đói, độ khát, vị trí và chiến lược hiện tại.

\section{Kiểm thử}

Project có sử dụng JUnit Jupiter để kiểm thử một số hành vi quan trọng:

\begin{itemize}
    \item \texttt{HunterStrategyTest}: kiểm tra chiến lược săn mồi di chuyển đến nguồn thức ăn gần nhất khi không có con mồi, đồng thời giữ đúng target khi đuổi con mồi.
    \item \texttt{TigerTest}: kiểm tra sát thương phục kích của hổ khi tấn công trong rừng và kiểm tra việc thỏ có thể đi vào bụi rậm trong khi sói không thể.
    \item \texttt{EntityManagerTest}: kiểm tra hươu đổi sang \texttt{AggressiveStrategy} khi độ đói xuống dưới ngưỡng nguy hiểm.
\end{itemize}

Các test này tập trung vào những hành vi quan trọng nhất của mô phỏng: chọn mục tiêu, tấn công, thay đổi chiến lược và ràng buộc di chuyển theo địa hình. Nhờ có test, nhóm có thể phát hiện lỗi khi sửa logic strategy hoặc thay đổi cách xử lý entity.

\section{Kết luận}

Nhóm đã xây dựng được ứng dụng mô phỏng hệ sinh thái hoang dã bằng JavaFX theo hướng lập trình hướng đối tượng. Chương trình thể hiện rõ các nguyên lý OOP thông qua hệ thống lớp kế thừa, interface, đa hình, đóng gói và trừu tượng hóa. Ngoài ra, project còn áp dụng Strategy Pattern để tách hành vi động vật, Factory để tạo con non và tách biệt logic mô phỏng với giao diện.

\subsection{Hạn chế}

\begin{itemize}
    \item Ảnh demo cần được bổ sung đầy đủ sau khi chạy chương trình thực tế.
    \item Một số loài sinh vật như cá và vịt có thể cần thêm sprite riêng để giao diện trực quan hơn.
    \item Test mới tập trung vào một số hành vi chính, chưa bao phủ toàn bộ vòng đời mô phỏng.
    \item Mô phỏng vẫn có thể mở rộng thêm nhiều quy luật sinh thái phức tạp hơn.
\end{itemize}

\subsection{Hướng phát triển}

\begin{itemize}
    \item Bổ sung nhiều loài động vật, thực vật và địa hình mới.
    \item Thêm màn hình cấu hình trước khi bắt đầu mô phỏng.
    \item Cải thiện giao diện thống kê để theo dõi biến động quần thể theo thời gian.
    \item Mở rộng test cho \texttt{WorldMap}, \texttt{SimulationEngine}, sinh sản và xử lý va chạm.
    \item Bổ sung thêm hiệu ứng hình ảnh và âm thanh để mô phỏng sinh động hơn.
\end{itemize}

\end{document}
